\ficha{Fritsch e Franco (1992), Aspects of the Brazilian Experience under the Gold Standard}{fritschfranco1992}

\begin{identificacao}
FRITSCH, Winston; FRANCO, Gustavo H. B. \textbf{Aspects of the Brazilian Experience under the Gold Standard}. Rio de Janeiro: Departamento de Economia, PUC-Rio, 1992. (Texto para Discussão, n. 286).

Texto para discussão, 23 páginas, lido integralmente, em inglês. Versão revista de trabalho apresentado na conferência \textit{The Gold Standard in the Countries of the Periphery}, Lisboa, 16 e 17 de dezembro de 1991. As páginas indicadas abaixo são as impressas no cabeçalho do original, que antecedem em duas a paginação do arquivo PDF.
\end{identificacao}

\bloco{Problema e tese}

O ensaio pergunta duas coisas ao mesmo tempo: que efeito o padrão-ouro produz numa economia primário-exportadora dependente de importação de capital, e se o mecanismo internacional de ajuste da Pax Britannica embutia viés sistêmico contra a periferia. A resposta é que o Brasil adotou o padrão-ouro nas duas ocasiões, 1906-1914 e 1926-1930, não como âncora contra a inflação mas como anteparo contra a apreciação cambial produzida pela conjunção de boa receita de exportação com auge de empréstimo internacional (p. 22). Como o desenho institucional escolhido era passivo, o ajuste externo transmitiu integralmente seu ciclo à base monetária, expandindo crédito na alta e contraindo com violência na baixa (p. 22).

\bloco{Argumento}

\begin{enumerate}[leftmargin=*,nosep]
\item O impasse imperial entre metalistas e papelistas era insolúvel enquanto houvesse ágio do ouro sobre o par legal de 1847, de 27 pence por mil-réis (p. 5). Foi a bonança do balanço de pagamentos, não o esforço deflacionista, que levou o câmbio a 27 pence em outubro de 1888 e tornou possível o compromisso, já que os metalistas não se opunham a emissão integralmente lastreada em ouro (p. 5). O Banco Nacional do Brasil, autorizado em 1889, ficou identificado com a Monarquia, e Rui Barbosa, papelista convicto, desfez o arranjo em 1890 (p. 6).

\item A depreciação dos anos 1890 deteriorou as finanças públicas e criou um problema de transferência interna, porque a receita em ouro era deficitária e a receita em papel superavitária (p. 7, nota 2). A acomodação financeira em Londres cobrou seu preço em política doméstica, com deflação severa, crise bancária e recessão na virada do século (p. 7-8).

\item Nos primeiros anos do novo século o quadro externo se inverteu: borracha em alta, retorno do Brasil ao mercado internacional de capitais, dívida externa federal subindo de 42 para 144 milhões de libras entre 1900 e 1913 e investimento direto de 105 para 295 milhões entre 1902 e 1913, com participação nos fluxos internacionais de investimento cinco vezes maior que a participação no comércio (p. 8). O câmbio saiu de menos de 7 pence em 1898 para 17,5 pence em agosto de 1905, e produtores de exportáveis e de substitutos de importação passaram a pressionar por estabilização cambial e pelo fim do dinheiro caro (p. 8).

\item Em 1906 o temor de que a compra da safra grande apreciasse ainda mais o câmbio levou à adoção do padrão-ouro com a criação da Caixa de Conversão (p. 9). A questão política era até onde deixar a apreciação ir: os produtores queriam fixar logo, os ortodoxos queriam caminhar de volta aos 27 pence, o que significaria apreciar quase 100\% sobre a taxa vigente (p. 9). O compromisso em 15 pence foi viabilizado pelo limite de emissão, exigido pelos deflacionistas justamente porque, mantida a bonança, o limite seria atingido e a taxa teria de ser revista para cima (p. 9). O limite, nessa leitura, não é teto prudencial e sim gatilho de revalorização.

\item Dois traços definem o arranjo. Era padrão-ouro na margem, pois só as notas da Caixa tinham lastro oficial em ouro (p. 9). E preservava o Banco do Brasil, tal como recriado em 1901, acumulando as funções de banco comercial e de banqueiro do governo, o que os autores leem como vindicação tardia do projeto papelista de grande banco nacional com traços de banco central e papel desenvolvimentista (p. 9). Na prática o Banco do Brasil podia manter a taxa dentro dos pontos do ouro por operações cambiais, mas não tinha ativos domésticos líquidos para neutralizar superávits externos grandes, e não havia disposição para esterilizar, dado o viés anti-inflacionário da opinião corrente, os déficits orçamentários e o tamanho reduzido do mercado de capitais interno (p. 9-10).

\item Daí o padrão-ouro limpo, em que os movimentos de ouro atingem inteiramente as condições domésticas de crédito e amplificam o ciclo externo (p. 10). A Tabela 3 mostra a emissão da Caixa subindo de 38 mil contos em 1906 para 406 em 1912, caindo a 295 em 1913 e a 158 em 1914, enquanto a base do Tesouro encolhia de 665 para 601 e só voltava a crescer com a guerra (p. 11). A expansão de 1906 a 1912 veio da borracha até 1910 e da conta de capital até 1912, com interrupção breve pela crise Knickerbocker de 1907-08; a reversão de 1913, com queda de preços de borracha e café e retração do capital após as Guerras Balcânicas, virou contração monetária prolongada, e a conversibilidade foi abandonada em 1914 para preservar o que restava do depósito da Caixa como reserva de guerra (p. 11).

\item O episódio de 1926-1930 repete o padrão, com a Caixa de Estabilização moldada exatamente sobre a Caixa de Conversão e taxa de conversão perto de 6 pence, 40\% da paridade do pré-guerra (p. 12).

\item A racionalidade política é a mesma nas duas vezes. Os dois episódios foram precedidos por colapsos cambiais e por negociações com banqueiros internacionais que impuseram as políticas restritivas de Campos Sales e de Artur Bernardes (p. 13). O padrão-ouro tinha apoio dos negócios domésticos, porque estancava a apreciação e acabava com o dinheiro caro, e das autoridades públicas de todos os níveis, porque restaurava a reputação de crédito e abria financiamento de longo prazo para obra pública e reescalonamento de dívida (p. 14). O ponto decisivo para o capítulo é que os ortodoxos não se opunham à adoção: sua objeção se concentrava no automatismo, no caráter não discricionário da gestão monetária, o que determinou o desenho da Caixa como órgão emissor passivo de lastro integral em vez de autoridade de reserva fracionária, e no nível em que fixar a taxa (p. 14).

\item Uma vez adotado, o padrão tinha valor político próprio, e por isso os governos sustentaram paridade sobrevalorizada até o último minuto nas duas crises, ao custo de contração monetária e perda de produto (p. 14). O padrão-ouro, contudo, não funcionou como dispositivo estabilizador: a inflação já havia sido debelada antes por política fiscal e monetária severa, e foi a melhora do balanço de pagamentos que criou a condição para recorrer ao ouro (p. 15).

\item Na parte sistêmica, os autores examinam três linhas de assimetria. A primeira, de Morgenstern e Lindert, atribui ao Banco da Inglaterra o poder de deslocar o ônus do ajuste para as extremidades do sistema elevando a taxa de redesconto (p. 16); Ford e Moggridge são céticos e exigem exame de casos (p. 16-17). A segunda, de Ford, aponta que cortar empréstimo é mais fácil que aumentar tomada de empréstimo, de modo que o credor ajusta melhor que o devedor (p. 17), e que ao devedor restam a contração dolorosa ou a saída do ouro (p. 18); permanecer no ouro exigia acumular reservas, o investimento improdutivo de que fala Furtado, e pagar senhoriagem aos emissores de moeda de reserva (p. 18). A terceira, de Triffin e Furtado, supõe padrão perverso das exportações de capital do centro, contestado por Cairncross e Rostow (p. 18-19).

\item A evidência brasileira não fecha nenhuma das três. A correlação entre movimentos de capital e termos de troca é de -0,02 em 1870-1900, de 0,304 no século XX e de 0,688, significativa a 1\%, em 1900-1918 (p. 20). Os autores concluem que não há padrão sistemático claro (p. 21), mas que um regime de correlação positiva se instala no século XX por efeito da valorização do café: a partir do Convênio de Taubaté, o financiamento externo dos estoques amarra o preço do café, e portanto os termos de troca brasileiros, à disponibilidade de crédito externo (p. 22). Na conclusão, a única assimetria que sobrevive é a de custos de ajuste entre credor e devedor (p. 23).
\end{enumerate}

\chave{The price eventually paid to obtain foreign financial accomodation to help stabilise the economy was a swing of domestic policies towards an extremely orthodox orientation, imposed by Rotschilds, the Brazilian goverment London bankers, as a condition for a large funding loan.}{p. 7}

\chave{This limit was explicitly demanded by deflationists based on the pressumption that if the balance of payments bonanza continued the limit would be reached and the exchange rate would have to be revised upwards.}{p. 9}

\chave{Orthodox policy-makers and rentiers tried to resist the abandonment of the ongoing austere monetary policies. However, they did not oppose the adoption of the gold standard.}{p. 14}

\bloco{Conceitos e definições operacionais}

Padrão-ouro na margem: regime em que apenas a emissão nova de um órgão específico tem lastro metálico, convivendo com o estoque prévio de papel inconversível do Tesouro (p. 9). Padrão-ouro limpo: aquele em que não há esterilização, e os movimentos de ouro incidem integralmente sobre as condições domésticas de crédito (p. 10). Automatismo: caráter não discricionário da gestão monetária, operacionalizado como emissão de lastro integral, um por um, em oposição a autoridade emissora de reserva fracionária (p. 14). Anteparo, e não âncora: a fixação cambial é adotada para deter apreciação, não para debelar inflação já controlada por outros meios (p. 15, p. 22). Transferência interna: descasamento entre receita e despesa públicas em ouro e em papel, que faz a desvalorização real piorar o saldo fiscal em moeda local (p. 7, nota 2).

\bloco{Evidência e método}

Revisão narrativa apoiada em séries agregadas, sem modelo formal nem econometria além de correlações bivariadas. Quatro tabelas, de investimento britânico e contas brasileiras em 1885-1892, de base monetária e câmbio em 1888-1899, e de câmbio real, base monetária, produto industrial e investimento em 1903-1917 e 1919-1930, além de dois gráficos de entrada de capital contra termos de troca, para 1870-1900 e 1901-1929 (p. 6, 7, 11, 12, 20, 21). Fontes secundárias: Goldsmith (1986), IBGE (1987), Mitchell (1978), Stone, Franco (1991) e Fritsch (1988). O contrafactual é implícito e comparativo, com o episódio de 1926-1930 servindo de réplica ao de 1906-1914.

\bloco{Agentes e vertentes}

Do lado brasileiro: metalistas e papelistas como campos do debate imperial, Mauá, Rui Barbosa como papelista, Campos Sales e Artur Bernardes como executores de ajuste ortodoxo, os Rothschild como banqueiros londrinos do governo brasileiro e impositores da condicionalidade do funding loan, o Banco Nacional do Brasil ligado à Banque de Paris et des Pays Bas, e o Banco do Brasil recriado em 1901. Nenhum agente individual de 1906 é nomeado: os autores tratam ortodoxos e produtores domésticos como blocos. A conclusão registra que metalistas e papelistas reapareceriam depois como desenvolvimentistas, estruturalistas, monetaristas, heterodoxos e ortodoxos (p. 22). Do lado da literatura: Morgenstern, Lindert, Ford, Moggridge, Eichengreen, Triffin, Furtado, Prebisch, Cairncross, Rostow, Saul e Gunder-Frank.

\bloco{Críticas e limites}

O texto é working paper, não artigo publicado, com erros tipográficos e de grafia no original. A tese sobre o limite de emissão, que é o achado mais útil para a dissertação, é afirmada sem referência documental: não se diz onde os deflacionistas o exigiram nem com que palavras. Os campos de 1906 aparecem como blocos sem nomes, sem imprensa e sem debate parlamentar, que é exatamente a lacuna que a dissertação preenche. Da cronologia da Caixa faltam a alteração da taxa para 16 pence e a ampliação do limite no período 1910-1913, a custódia do ouro em 1911 e a agência de Londres; o papel do Banco do Brasil é afirmado como possibilidade, não documentado em operação. As correlações são bivariadas e sem controle, limitação que os próprios autores apontam ao criticar Gonçalves e Barros (p. 20). A conclusão sobre assimetria sistêmica é fraca por construção, já que o próprio texto conclui pela ausência de padrão regular (p. 21).

\begin{dialogo}
\item Procurar nos jornais e nos Documentos Parlamentares o argumento de que o limite de emissão da Caixa era provisório e serviria para forçar a revisão da taxa para cima. Vocabulário: limite de emissão, elevação da taxa, par legal, 27 pence, restauração do valor. Confirma ou nega a tese de que o limite foi exigência deflacionista, e não prudência técnica (p. 9).
\item Testar se os ortodoxos atacaram a Caixa como instituição ou apenas seu automatismo e sua taxa. Vocabulário: emissão lastreada, não é papel-moeda, conversibilidade, lastro integral, discrição do Tesouro. Se a imprensa ortodoxa aceitar a Caixa e brigar só pela taxa, confirma a afirmação da p. 14 e reforça o eixo ortodoxo contra expansionista no lugar de metalista contra papelista.
\item Procurar as propostas de dar ao Banco do Brasil funções cambiais ou emissoras concorrentes com a Caixa. Vocabulário: carteira cambial, banco de emissão, banco nacional, pontos do ouro. Testa a leitura de que o arranjo de 1906 foi vindicação tardia do projeto papelista de grande banco (p. 9), e mapeia o conflito de competência entre Caixa, Tesouro e Banco do Brasil.
\item Confrontar a tese pró-cíclica com a percepção de época: procurar queixas de abundância de numerário e de especulação entre 1909 e 1912 e de aperto e crise de crédito em 1913 e 1914, ao lado dos boletins diários de movimento da Caixa. Testa se os contemporâneos leram a expansão e a contração como efeito do ouro entrando e saindo da Caixa (p. 10-11).
\item Procurar a associação explícita entre empréstimos de valorização do café e entrada de ouro na Caixa, e entre a Caixa e a restauração do crédito externo para obras públicas. Vocabulário: crédito no estrangeiro, confiança, empréstimo de valorização, Taubaté. Testa a tese da p. 14 sobre a motivação de reputação de crédito e a da p. 22 sobre a valorização como divisor de águas.
\end{dialogo}

\usonocapitulo{Entra na parte II, sobre banqueiros e enforcement, como o elo brasileiro explícito: o funding loan imposto pelos Rothschild como condicionalidade (p. 7) e a adoção do padrão-ouro em 1906 como instrumento de restauração da reputação de crédito, que abre financiamento de longo prazo para obra pública e reescalonamento (p. 14). Sustenta a afirmação de que a ortodoxia de 1906 não era contrária à Caixa, e que a disputa se deu sobre automatismo e nível da taxa, o que ampara o eixo analítico ortodoxo contra expansionista. Faz ponte com a parte I ao qualificar a assimetria centro-periferia: a única assimetria que o texto sustenta com evidência brasileira é a de custos de ajuste entre credor e devedor (p. 23), e não a transferência de ônus pela taxa do Banco da Inglaterra. Serve ainda à parte IV como descrição do desenho institucional da Caixa, órgão emissor passivo de lastro integral com limite legal (p. 9, p. 14).}
